% ---------- FRONT MATTER: ROMAN NUMBERS ----------
\pagenumbering{roman}
\setcounter{page}{1}

\begin{titlepage}
  \thispagestyle{empty}
  \centering
  {\Huge \textbf{ScamGuard: AI-Based Online Scam}}\\[0.25cm]{\Huge \textbf{and Fraud Detection System}}\\[0.8cm]
  \includegraphics[width=0.18\linewidth]{VTU-Logojpg.jpg}\\[0.7cm]
  {\Large \textbf{PROJECT REPORT}}\\[0.15cm]
  {\large Major Project Phase-2 (23CS7PWPP2), Semester VII}\\[0.3cm]
  {\large \textit{Submitted by}}\\[0.5cm]
  {\large Bhavya Goyal (1BM23CS063)}\\
  {\large Priyam Thakur (1BM23CS252)}\\
  {\large Samir Chaudhary (1BM23CS294)}\\
  {\large Sana Subodh (1BM23CS296)}\\[0.7cm]
  {\itshape in partial fulfillment for the award of the degree of}\\[0.15cm]
  {\Large \textbf{BACHELOR OF ENGINEERING}}\\[0.2cm]
  {\large in}\\[0.2cm]
  {\Large \textbf{COMPUTER SCIENCE AND ENGINEERING}}\\[0.9cm]
  {\large Under the Guidance of}\\
  {\large \textbf{Dr. Kayarvizhy N}}\\
  {\large Professor, Department of CSE, BMSCE}
  \vfill
  \includegraphics[width=0.18\linewidth]{bmscelogo.jpg}\\[0.5cm]
  {\Large \textbf{B.M.S. COLLEGE OF ENGINEERING}}\\
  {\large (Autonomous Institution under VTU)}\\
  {\large BENGALURU - 560019}\\
  {\large Academic Year - 2026--2027}\\[0.2cm]
  {\large Date of Submission: November 2026}
\end{titlepage}

\chapter*{Certificate}
\thispagestyle{empty}
\vspace{2cm}
This is to certify that the project report entitled \textbf{``ScamGuard: AI-Based Online Scam and Fraud Detection System''} is a bonafide work carried out by \textbf{Bhavya Goyal (USN: 1BM23CS063)}, \textbf{Priyam Thakur (USN: 1BM23CS252)}, \textbf{Samir Chaudhary (USN: 1BM23CS294)} and \textbf{Sana Subodh (USN: 1BM23CS296)} as Major Project Phase-2 (Course Code: 23CS7PWPP2) in the seventh semester, in partial fulfillment of the requirements for the award of the degree of \textbf{Bachelor of Engineering} in \textbf{Computer Science and Engineering}, under my supervision during the academic year 2026--2027.

\vspace{2cm}
\begin{flushleft}
Guide Signature\hspace{30mm} HoD signature \hspace{25mm} Principal Signature\\
(Dr. Kayarvizhy N) \hspace{22mm} Dr. Kavitha Sooda \hspace{19mm} Dr. Bheemsha Arya
\end{flushleft}
\vspace{2cm}
\begin{flushleft}
Name and signature of Examiners along with date \\
\vspace{1cm}
Examiner 1 \\
\vspace{1cm}
Examiner 2 \\
\end{flushleft}

\chapter*{Abstract}
\thispagestyle{empty}
Fraudulent and unwanted SMS messages are one of the cheapest ways to reach a phone user in India. They imitate bank KYC reminders, electricity disconnection notices, parcel updates and job offers, and they increasingly arrive in Hindi or in a mix of Hindi and English. ScamGuard is an AI-based system that classifies a message as scam or genuine, explains the verdict at phrase level and tells the user what to do next.

The core detector is trained on ScamGuard-IN-2026, a corpus of 11,399 Indian messages built from user-reported smishing, real Indian SMS, 2025--26 advisory texts and synthetic messages used for training only. Reported messages were audited: 19.5\% of retained reports were genuine bank alerts, OTP messages or advertisements. Near-duplicate campaigns were grouped with MinHash and kept on one side of the train--test split. A logistic-regression classifier over word and character TF-IDF features, combined with 31 interpretable intent patterns, reaches 97.1\% accuracy with a 3.2\% false-alarm rate on 1,874 held-out real messages from unseen campaigns; training on unverified report labels raises the false-alarm rate to 34.0\%.

As the advancement suggested by our guide, we studied cross-domain generalization and lightweight domain adaptation. Multinomial Naive Bayes and Logistic Regression trained on the UCI SMS Spam Collection (UK) were tested, without retraining, on an Indian SMS dataset. Logistic Regression fell from a spam F1-score of 0.911 to 0.678 and missed about half of the Indian spam. Adding 10\% of a labelled Indian adaptation pool (143 messages) raised the mean target F1 over three seeds from 0.681 to 0.964, without lowering source-domain performance; the gain held in every seed, was significant by McNemar's test, and persisted when near-duplicate messages were kept apart. The system is deployed as a web application (Next.js front end, FastAPI back end, PostgreSQL), with a dedicated cross-domain page and a Streamlit application, and every number in this report is reproducible from the project repository.

\vspace{0.4cm}\noindent\textbf{Keywords:} scam detection, SMS spam, smishing, domain shift, domain adaptation, TF-IDF, logistic regression, Naive Bayes, label noise, data leakage

\chapter*{DECLARATION}
\thispagestyle{empty}
We hereby declare that the Major Project Phase-2 work entitled ``ScamGuard: AI-Based Online Scam and Fraud Detection System'' is a bonafide work and has been carried out by us under the guidance of Dr. Kayarvizhy N, Professor, Department of Computer Science and Engineering, B.M.S. College of Engineering, Bengaluru, in partial fulfillment of the requirements of the degree of Bachelor of Engineering in Computer Science and Engineering of Visvesvaraya Technological University, Belagavi.
We further declare that, to the best of our knowledge and belief, this project has not been submitted either in part or in full to any other university for the award of any degree.
\vspace{0.3cm}
\begin{flushleft}
Candidate details:
\end{flushleft}
\vspace{0.1cm}
\begin{center}
\begin{tabular}{|c|l|c|c|}
\hline
\textbf{SL. NO.} & \textbf{Student Name} & \textbf{USN} & \textbf{Student's Signature} \\
\hline
1 & Bhavya Goyal & 1BM23CS063 & \\
\hline
2 & Priyam Thakur & 1BM23CS252 & \\
\hline
3 & Samir Chaudhary & 1BM23CS294 & \\
\hline
4 & Sana Subodh & 1BM23CS296 & \\
\hline
\end{tabular}
\end{center}
\begin{flushleft}
Place: Bengaluru \\
Date: \\
\end{flushleft}
Certified that these candidates are students of Computer Science and Engineering Department of B.M.S. College of Engineering. They have carried out the project work titled ``ScamGuard: AI-Based Online Scam and Fraud Detection System'' as Major Project Phase-2 work. It is in partial fulfillment for completing the requirement for the award of B.E. degree by VTU. The work is original and we duly certify the same.
\vspace{1cm}
\begin{flushleft}
Guide Name: Dr. Kayarvizhy N\\
Date: \\
Signature \\
\end{flushleft}

\chapter*{Acknowledgment}
\thispagestyle{empty}
We are grateful to our guide, Dr. Kayarvizhy N, Professor, Department of Computer Science and Engineering, for steering this project towards questions that could be answered with evidence: first, a dataset that reflects the messages people in India actually receive, and then the cross-domain generalization and adaptation study that she proposed as the advancement of this work.

We thank Dr. Kavitha Sooda, Head of the Department of Computer Science and Engineering, and Dr. Bheemsha Arya, Principal, B.M.S. College of Engineering, for the facilities and the academic environment that made this work possible. We also thank the faculty members who reviewed our phase presentations and pointed out gaps in our evaluation.

We acknowledge the researchers who released the IMC 2025 smishing collection, the UCI SMS Spam Collection and the India Spam SMS collection under open licences. Finally, we thank our families and friends, several of whom forwarded the suspicious messages that shaped our test cases.
\newpage
\tableofcontents
\thispagestyle{empty}
\newpage
\listoffigures
\thispagestyle{empty}
\newpage
\listoftables
\thispagestyle{empty}

\clearpage
\pagenumbering{arabic}
\setcounter{page}{1}
